\documentclass[12pt,a4paper]{article}
\usepackage[utf8]{inputenc}
\usepackage[T1]{fontenc}
\usepackage{amsmath,amssymb,amsthm}
\usepackage{geometry}
\usepackage{enumitem}
\usepackage{hyperref}
\usepackage{fancyhdr}
\usepackage{titlesec}

\geometry{margin=1in}
\pagestyle{fancy}
\fancyhf{}
\rhead{Lecture 9}
\lhead{Quantum Systems via Linear Algebra}
\cfoot{\thepage}

\titleformat{\section}{\large\bfseries}{}{0em}{}
\titleformat{\subsection}{\normalsize\bfseries}{}{0em}{}

\newtheorem{definition}{Definition}
\newtheorem{theorem}{Theorem}

\title{\textbf{Lecture 9: Continuous Systems and the Fourier Transform} \\ \large $L^2(\mathbb{R})$, Position and Momentum Bases, and the Canonical Commutator}
\author{Quantum Systems via Linear Algebra, Differential Equations, and Probability}
\date{}

\begin{document}
\maketitle
\thispagestyle{fancy}

% ============================================================
\section{1. The State Space $L^2(\mathbb{R})$}
% ============================================================

Up to now, all state spaces have been finite-dimensional ($\mathbb{C}^2$, $\mathbb{C}^4$, etc.).  We now move to systems where the state space is the \emph{infinite-dimensional} Hilbert space $L^2(\mathbb{R})$---the space of square-integrable complex-valued functions on the real line.

\begin{definition}[$L^2(\mathbb{R})$]
The set of all functions $v : \mathbb{R} \to \mathbb{C}$ satisfying
\[
    \int_{-\infty}^{\infty} |v(x)|^2\,dx < \infty.
\]
This is a complex vector space with inner product
\[
    \langle u, v \rangle = \int_{-\infty}^{\infty} u^*(x)\,v(x)\,dx.
\]
\end{definition}

\noindent\textbf{Engineering connection:} $L^2(\mathbb{R})$ is exactly the space of finite-energy signals from signal processing.  The inner product is the cross-correlation at zero lag.  The norm $\|v\| = \sqrt{\langle v, v \rangle}$ is the signal energy.

\subsection{1.1 Transition from finite to infinite dimensions}
The passage from $\mathbb{C}^N$ to $L^2(\mathbb{R})$ is notational:
\begin{itemize}[nosep]
    \item \textbf{Vectors}: column of $N$ entries $\to$ function $v(x)$.
    \item \textbf{Inner product}: $\sum_i u_i^* v_i \to \int u^*(x)\,v(x)\,dx$.
    \item \textbf{Kronecker delta}: $\delta_{ij} \to \delta(x-x')$ (Dirac delta).
    \item \textbf{Basis expansion}: $\sum_i \to \int$.
\end{itemize}

\subsection{1.2 Physical interpretation}
For a particle on a line:
\begin{itemize}[nosep]
    \item $v(x)$ is the \textbf{position representation}: $|v(x)|^2\,dx$ is the probability of finding the particle between $x$ and $x+dx$.
    \item $\tilde{v}(p)$ (Fourier transform of $v$) is the \textbf{momentum representation}: $|\tilde{v}(p)|^2\,dp$ is the probability of the momentum lying between $p$ and $p+dp$.
\end{itemize}

% ============================================================
\section{2. The Fourier Transform as a Change of Basis}
% ============================================================

\subsection{2.1 Complex exponentials}
For a fixed number $p$, the function $e^{ipx} = \cos(px) + i\sin(px)$ is a sinusoid with frequency $p/(2\pi)$.

\subsection{2.2 Fourier transform pair}
Any $v \in L^2(\mathbb{R})$ can be decomposed into complex exponentials:
\[
    \boxed{v(x) = \frac{1}{\sqrt{2\pi}}\int_{-\infty}^{\infty} \tilde{v}(p)\,e^{ipx}\,dp,}
\]
where $\tilde{v}(p)$ is the \textbf{Fourier transform}:
\[
    \boxed{\tilde{v}(p) = \frac{1}{\sqrt{2\pi}}\int_{-\infty}^{\infty} v(x)\,e^{-ipx}\,dx.}
\]
$\tilde{v}(p)$ gives the amplitude of the frequency-$p$ component in the signal $v(x)$.

\medskip
\noindent\textbf{Engineering connection:} This is the standard Fourier transform from signal processing.  The position representation $v(x)$ is the ``time-domain'' signal; the momentum representation $\tilde{v}(p)$ is the ``frequency-domain'' spectrum.  Parseval's theorem $\int|v(x)|^2 dx = \int|\tilde{v}(p)|^2 dp$ ensures both representations carry the same total probability (energy).

% ============================================================
\section{3. The Dirac Delta Function}
% ============================================================

\begin{definition}[Delta function]
The Dirac delta function $\delta(x-x')$ is defined operationally by its action on any reasonable function $f$:
\[
    \int_{-\infty}^{\infty} \delta(x-x')\,f(x)\,dx = f(x').
\]
\end{definition}

\subsection{3.1 Key properties}
\begin{enumerate}[nosep]
    \item $\delta(x-x')\,f(x) = f(x')\,\delta(x-x')$ \quad (only nonzero when $x=x'$).
    \item $\displaystyle\int_{-\infty}^{\infty}\delta(x-x')\,dx = 1$.
    \item Combining 1 and 2: $\displaystyle\int \delta(x-x')\,f(x)\,dx = f(x')$.
\end{enumerate}

\subsection{3.2 Fourier representation of the delta function}
Substitute the Fourier transform $\tilde{v}(p)$ into the inverse transform for $v(x')$:
\[
    v(x') = \frac{1}{\sqrt{2\pi}}\int \tilde{v}(p)\,e^{ipx'}dp
    = \frac{1}{\sqrt{2\pi}}\int\left(\frac{1}{\sqrt{2\pi}}\int v(x)\,e^{-ipx}dx\right)e^{ipx'}dp.
\]
Interchange the order of integration:
\[
    v(x') = \int v(x)\underbrace{\left(\frac{1}{2\pi}\int e^{ip(x'-x)}dp\right)}_{\text{must be }\delta(x'-x)}dx.
\]
This yields the fundamental identity:
\[
    \boxed{\delta(x-x') = \frac{1}{2\pi}\int_{-\infty}^{\infty} e^{ip(x-x')}\,dp.}
\]
By symmetry (interchanging $x\leftrightarrow p$):
\[
    \delta(p-p') = \frac{1}{2\pi}\int_{-\infty}^{\infty} e^{ix(p-p')}\,dx.
\]

% ============================================================
\section{4. Position and Momentum as Conjugate Bases}
% ============================================================

\subsection{4.1 Position basis}
The position observable $\hat{x}$ acts on functions by multiplication: $\hat{x}\,v(x)=x\,v(x)$.  Its ``eigenfunctions'' are delta functions: $\hat{x}\,\delta(x-x_0) = x_0\,\delta(x-x_0)$.  They satisfy the continuous orthogonality relation $\int \delta(x-a)\,\delta(x-b)\,dx = \delta(a-b)$.

Any state can be expanded:
\[
    v(x) = \int v(x')\,\delta(x-x')\,dx'.
\]

\subsection{4.2 Momentum basis}
The momentum observable $\hat{p}$ has eigenfunctions $e^{ipx}/\sqrt{2\pi}$ (plane waves, completely delocalized in position).  The Fourier transform pair:
\[
    v(x) = \frac{1}{\sqrt{2\pi}}\int \tilde{v}(p)\,e^{ipx}\,dp, \qquad
    \tilde{v}(p) = \frac{1}{\sqrt{2\pi}}\int v(x)\,e^{-ipx}\,dx.
\]
The Fourier transform \emph{is} the change-of-basis matrix from position to momentum representation.

% ============================================================
\section{5. The Momentum Operator in Position Space}
% ============================================================

Acting on a position-space function, the momentum operator is represented as a derivative:
\[
    \boxed{\hat{p}\,v(x) = -i\frac{d}{dx}v(x).}
\]

\subsection{5.1 Verification}
A plane wave $v(x)=e^{ipx}$ is a momentum eigenfunction:
\[
    \hat{p}\,e^{ipx} = -i\frac{d}{dx}e^{ipx} = -i(ip)\,e^{ipx} = p\,e^{ipx}. \quad\checkmark
\]

\subsection{5.2 Canonical commutation relation}
Compute $[\hat{x},\hat{p}]$ acting on an arbitrary $v(x)$:
\begin{align}
    \hat{x}\hat{p}\,v &= x\bigl(-iv'(x)\bigr) = -ixv'(x), \\
    \hat{p}\hat{x}\,v &= -i\frac{d}{dx}\bigl(xv(x)\bigr) = -i\bigl(v(x)+xv'(x)\bigr).
\end{align}
Subtracting: $(\hat{x}\hat{p}-\hat{p}\hat{x})v = -ixv'+iv+ixv' = iv$.  Since this holds for all $v$:
\[
    \boxed{[\hat{x},\hat{p}] = i} \qquad (\text{with } \hbar=1).
\]
This is the \textbf{canonical commutation relation}.  Since $[\hat{x},\hat{p}]\neq 0$, position and momentum \textbf{cannot be simultaneously measured}---the basis of the uncertainty principle (Lecture~10).

\noindent\textbf{Restoring $\hbar$:} $[\hat{x},\hat{p}]=i\hbar$.

% ============================================================
\section{6. Simple System Matrices for Particles}
% ============================================================

\subsection{6.1 Constant-velocity propagation: $H=cP$}
The Schr\"odinger equation $i\frac{\partial v}{\partial t}=c(-i\frac{\partial v}{\partial x})$ simplifies to:
\[
    \frac{\partial v}{\partial t} + c\frac{\partial v}{\partial x} = 0.
\]
General solution: $v(x,t)=f(x-ct)$ for any shape $f$.  The entire probability distribution moves rigidly to the right at speed $c$ without changing shape.

\medskip
\noindent\textbf{Engineering analog:} This is a pure delay line---the signal propagates without dispersion.

\subsection{6.2 Free particle: $H=\hat{p}^2/2m$}
The standard Schr\"odinger equation for a free particle of mass $m$:
\[
    \boxed{i\hbar\frac{\partial v}{\partial t} = -\frac{\hbar^2}{2m}\frac{\partial^2 v}{\partial x^2}.}
\]
Unlike $H=cP$, different frequency components now move at different speeds---the signal \textbf{disperses} (spreads out) over time.  The expectation value of position still obeys $\frac{d}{dt}\langle x \rangle=\frac{\langle p \rangle}{m}$, but the wave packet broadens.

\medskip
\noindent\textbf{Engineering analog:} This is a dispersive channel.  The group velocity $v_g = d\omega/dk = \hbar k/m$ depends on wavenumber, so narrow pulses spread.

% ============================================================
\section{7. Summary}
% ============================================================

\begin{enumerate}[nosep]
    \item The state space for continuous systems is $L^2(\mathbb{R})$---the space of finite-energy complex signals.
    \item The \textbf{Fourier transform} is the change-of-basis between position and momentum representations.
    \item The \textbf{Dirac delta function} $\delta(x-x')=\frac{1}{2\pi}\int e^{ip(x-x')}dp$ is the continuous orthogonality relation.
    \item Position eigenfunctions are delta functions; momentum eigenfunctions are plane waves $e^{ipx}/\sqrt{2\pi}$.
    \item The momentum operator is $\hat{p}=-id/dx$; $[\hat{x},\hat{p}]=i\hbar$.
    \item $H=cP$: dispersion-free propagation (delay line).
    \item $H=\hat{p}^2/2m$: dispersive propagation; wave packets broaden over time.
\end{enumerate}

\end{document}
